\section{场景颜色与后处理}\ruleline
看视口里的椅子，椅背有自己的颜色，椅腿会投影到地面，调高曝光后整幅画面又一起变亮。这几件事发生在不同的绘制阶段。开始时，\code{DrawModel} 把各子网格、位置、材质和实体编号记下来。\code{SceneRenderer} 按下图的次序使用它们。主画面把同材质的记录排在一起；共享网格和同组贴图的不透明物体也可用一次实例化绘制提交，仍能分别选中。

\begin{figure}[H]\centering
\begin{tikzpicture}[module/.append style={text width=5.6cm,minimum height=1.35cm,inner sep=6pt}]
\node[module] (a) at (0,0) {\textbf{1 画二维物体，收集三维物体}\\同时取出这一帧的灯光数据};
\node[module] (b) at (7.8,0) {\textbf{2 从灯光位置看场景}\\Shadow 记遮挡；RSM 记表面反射信息};
\node[module] (c) at (0,-2.25) {\textbf{3 从相机位置记录表面}\\深度记远近；法线记表面朝向};
\node[module] (d) at (7.8,-2.25) {\textbf{4 用辅助图估算光照}\\SSAO 算局部遮蔽；RSM 算间接光};
\node[module] (e) at (0,-4.5) {\textbf{5 混合本帧与前帧间接光}\\RSM History 减少采样结果的闪动};
\node[module] (f) at (7.8,-4.5) {\textbf{6 绘制模型的表面颜色}\\Renderer3D::EndScene 使用材质和光照};
\node[module] (g) at (0,-6.75) {\textbf{7 补上天空盒与调试线框}\\结果仍保存到场景的颜色图中};
\node[module] (h) at (7.8,-6.75) {\textbf{8 调整亮度，加入泛光}\\PostProcess 输出供编辑器显示的图像};
\draw[flow](a)--(b);\draw[flow](c)--(d);\draw[flow](e)--(f);\draw[flow](g)--(h);
\draw[flow](b.south)--++(0,-.35)-|(c.north);
\draw[flow](d.south)--++(0,-.35)-|(e.north);
\draw[flow](f.south)--++(0,-.35)-|(g.north);
\end{tikzpicture}\caption{当前主要顺序。是否执行某项效果，还取决于开关、灯光与所需图像是否可用。}\end{figure}
\subsection{阴影、深度与表面颜色}
先从方向光望向场景，记下光线遇到的最近表面，地面着色时便能判断椅腿是否挡住灯光。打开反射阴影图（RSM）时，还会记录灯所照见的表面颜色等信息，供后面估算墙面反射的光。换回主相机后，深度给出每个屏幕位置可见的表面，法线给出这个表面的朝向。屏幕空间环境遮蔽（SSAO）比较邻近表面，估计椅腿与地面接触处能收到多少环境光。

间接光计算把眼前的表面位置投回上一帧，用先前的颜色减少闪动。投到画面外、相机后方或裁剪范围外的点会跳过；切换场景或修改投影、光照及采样设置也会清掉旧图。运动物体和刚露出的表面仍可能留下拖影，因为当前没有保存物体的运动和上一帧对应的深度。同投影的运行时相机互换也尚不能单独识别。

\subsection{曝光、泛光与编辑器视口}
\code{ForgeLayer} 准备一个 \code{Framebuffer} 接收场景颜色、实体编号、观察空间法线和深度。颜色先写成浮点数，灯或天空的亮度可以超过 1。开启泛光时，\code{PostProcess::Render} 从这张高动态范围（HDR）图提取亮部，再把不同尺寸的亮部图像合成柔和的光晕。泛光默认用六层，可以设置为一到六层。最后，后处理按照曝光把亮度映射到显示范围，编辑器将最终图像交给 ImGui 的视口面板。

目前 \code{SceneRenderer} 持有一个 \code{PostProcess} 来管理曝光、泛光和输出图。其他效果仍有静态共享数据，两个视口的图像尚未全部分开。调试面板里能看到的深度预览，是同步读回整张深度图后拉伸得到的灰度。若要从灰度求真实距离，还要用写入深度时的相机投影还原。
